% Reporte técnico del Reto AgroCebada FIRA 2026. Formato IEEEtran con portada propia.
% Compilar desde reporte/entrega (dos veces, para las referencias cruzadas):
%   xelatex -interaction=nonstopmode main.tex && xelatex -interaction=nonstopmode main.tex
% Después, `python3 plataforma.py` genera el texto de cada campo de la plataforma y el conteo de palabras.
\documentclass[conference,letterpaper]{IEEEtran}
\usepackage{fontspec}
\setmainfont{TeX Gyre Termes}
\usepackage{polyglossia}
\setdefaultlanguage{spanish}
\usepackage{amsmath,amssymb}
\usepackage{graphicx}
\usepackage{booktabs}
\usepackage{array}
\usepackage{xcolor}
\usepackage{hyperref}
\hypersetup{colorlinks=true,linkcolor=black,citecolor=black,urlcolor=blue!50!black}
\graphicspath{{figuras/}}

\begin{document}

% Sin portada. Título, equipo y autores como encabezado de artículo, para ahorrar espacio.
\title{Rendimiento de cebada de temporal con régimen de registro y rezago espacial\\[2pt]
\large Reporte técnico. Reto AgroCebada FIRA 2026}
\author{\IEEEauthorblockN{Equipo Los Bien Dotados\\
José Alán Pomares Valdés, Luis Emmanuel Ramiro Luna y Vicente Josimar López Martínez\\
Asesor Luis Gustavo Olivares Soto}
\IEEEauthorblockA{Facultad de Ciencias de la Computación, Benemérita Universidad Autónoma de Puebla, Puebla, México\\
\href{mailto:pv202351433@alm.buap.mx}{pv202351433@alm.buap.mx}}}
\maketitle

\section{Introducción}\label{sec:intro}

El Reto AgroCebada FIRA 2026 pide estimar el rendimiento de grano de cebada de temporal, en toneladas por hectárea
(t/ha), de 59 parcelas de Hidalgo, Puebla y Tlaxcala~\cite{bases}. Estas entidades se cuentan entre las principales
productoras de cebada del país~\cite{sader2025}. FIRA busca que la estimación apoye decisiones productivas, de planeación
y eventualmente financieras~\cite{preguntas}.

Contamos con 197 parcelas de nueve municipios, todas del ciclo primavera-verano 2025. En 138 conocemos el rendimiento y
las usamos para entrenar. Las otras 59 son las que se califican, principalmente con la raíz del error cuadrático medio (RMSE), el error absoluto medio y el coeficiente de
determinación, sin ponderar por superficie~\cite{preguntas}. Para todas hay imágenes de los satélites
Sentinel-2~\cite{drusch2012}, Landsat y PlanetScope~\cite{frazier2021}, lluvia~\cite{funk2015},
temperatura~\cite{verdin2020} y topografía.

El rendimiento es el peso en báscula de toda la cosecha, sin ajuste por humedad ni impurezas. Lo reportó sobre todo la
comercializadora que compró el grano y, en menor medida, el productor~\cite{preguntas}. Esto importa porque los
rendimientos reportados suelen traer error~\cite{lobell2020}.

Predecir el rendimiento con satélites es un tema muy estudiado~\cite{vanklompenburg2020,weiss2020}. Aquí el problema
cambia por dos razones. La primera es que 138 parcelas son pocas, y la diferencia entre dos modelos suele ser menor que
el azar de la validación. La segunda es la forma en que se registró el rendimiento.

Al revisar las etiquetas encontramos que cada municipio registra el rendimiento de una sola de dos formas (anexo
técnico, Fig.~\ref{A-fig:regimenes}). En cuatro municipios los valores tienen dos decimales y promedian 4.49\,t/ha. Lo
llamamos régimen A. En los otros cinco están redondeados y promedian 3.05\,t/ha. Lo llamamos régimen B. La diferencia entre ambos, 1.4\,t/ha, es casi tres veces la variación dentro de cada uno,
así que conocer el régimen ya explica buena parte del rendimiento. No sabemos si
el rendimiento se midió distinto o solo se anotó distinto. FIRA pidió tratar todos los registros como una sola variable,
así que usamos el régimen como una característica del municipio y no como dos problemas separados.

Por eso nuestro trabajo no se limita a buscar el algoritmo con menor error. Buscamos responder qué información predice
de verdad el rendimiento, qué tan confiable es el error que estimamos y qué significa ese error para un productor o un
analista.

El resultado es un modelo sencillo. Usa el régimen del municipio y el rendimiento de las parcelas vecinas, y esperamos
que su RMSE en las 59 parcelas sea de 0.46 a 0.49\,t/ha. Lo acompañamos con un
intervalo por parcela, una traducción del error a pesos y una aplicación web.

\section{Justificación}\label{sec:justificacion}

FIRA pidió que la estimación sirva para decisiones productivas, de planeación y eventualmente financieras, y que cada
equipo explique qué significa su error en términos productivos~\cite{preguntas}. Para cualquiera de esas decisiones un
número solo no basta. Hace falta saber cuánto puede fallar.

Las mediciones de campo son escasas y costosas, mientras que las imágenes de satélite y el clima son públicos y cubren
cualquier parcela. Vale la pena saber cuánta información del rendimiento contienen realmente. Con pocas parcelas y un
registro heterogéneo, también vale la pena medir el error con cuidado, porque un error subestimado haría parecer
seguras decisiones que no lo son.

Este trabajo aporta una estimación por parcela con un error medido de forma conservadora, su traducción a pesos y una
aplicación para consultarla. El detalle de cada decisión está en los reportes técnicos del repositorio~\cite{repo}.

\section{Objetivos}\label{sec:objetivos}

\textbf{Objetivo general.} Estimar el rendimiento de cebada de temporal de las 59 parcelas del reto con el menor error
posible en parcelas no vistas, y medir con honestidad qué tan confiable es esa estimación.

\textbf{Objetivos específicos.}
\begin{enumerate}
\item Caracterizar las etiquetas de rendimiento y medir cómo su forma de registro limita lo que un modelo puede predecir.
\item Construir el modelo con una regla de decisión fijada antes de probar ideas, y conservar solo las piezas que
mejoren la validación de forma consistente.
\item Estimar el error esperado incluyendo el costo de elegir entre variantes, y entregar por parcela un intervalo de
predicción con cobertura verificada.
\item Responder con evidencia qué índice satelital, qué variable climática y qué combinación de variables pesan más.
\item Traducir el error a pesos por hectárea y a una señal de riesgo de crédito, e indicar en qué momento del ciclo puede
usarse la estimación.
\item Publicar una aplicación web que ejecute el modelo y muestre predicciones, intervalos y probabilidades.
\end{enumerate}

\section{Metodología}\label{sec:metodologia}

\textbf{Limpieza de las series.} Consideramos limpia una observación satelital con nubosidad de 20\,\% o menos y con
valor del índice. Elegimos ese umbral porque la nubosidad casi siempre vale 0 o 100\,\% (86\,\% de los registros) y
cuando vale 100 los índices vienen vacíos. Así cada parcela conserva de 19 a 34 observaciones de Sentinel-2 entre mayo
y octubre de 2025, con mediana de 25. Las observaciones descartadas no se imputan. Excluimos el índice de área foliar
(LAI) de Sentinel-2 porque su fórmula da valores negativos sobre suelo desnudo, sin significado físico, en 42\,\% de los
registros~\cite{ecuaciones}. Con las series limpias construimos 91 variables por parcela, como el máximo del NDVI de
2025, la fenología~\cite{zadoks1974}, la historia de 2022 a 2024, el clima y la topografía. Ninguna usa el rendimiento.

\textbf{Régimen.} Asignamos el régimen con una regla fija por nombre de municipio. No usamos el formato de cada
etiqueta, porque 7 de las 87 parcelas del régimen A tienen un solo decimal y esa regla las habría mandado al B. La regla
se conoce para las 59 parcelas, porque todas pertenecen a municipios con parcelas de entrenamiento. En seis parcelas de municipios con una a tres parcelas de entrenamiento el régimen se
dedujo de su propia etiqueta y de la cercanía a municipios del mismo formato. Sin ellas el RMSE pasa de 0.458 a
0.464\,t/ha, así que el efecto es pequeño.

\textbf{Esquema de evaluación.} Todo error se mide en parcelas que el modelo no vio. Usamos validación cruzada de 5
grupos, estratificada por municipio y repetida 10 veces, porque imita la partición oficial. Es el esquema A. Como la
validación al azar puede ser optimista cuando parcelas cercanas se parecen~\cite{roberts2017,ploton2020}, añadimos un
esquema B que deja fuera zonas completas, formadas al unir parcelas a menos de 2\,km. Todo lo que usa el rendimiento se
recalcula dentro de cada grupo con solo las parcelas de entrenamiento.

\textbf{Regla de decisión.} Con 138 parcelas una mejora puede aparecer por azar. Por eso una idea se acepta solo si baja
el RMSE en al menos 7 de las 10 repeticiones y si pasa una prueba de permutación con $p\le0.05$~\cite{ojala2010}. En la
prueba revolvemos las variables nuevas dentro de cada régimen y repetimos la validación 50 veces. Si el modelo predice
igual con ellas revueltas, la mejora no viene de esas variables.

\textbf{Rezago espacial.} Parcelas cercanas tienden a rendir parecido porque comparten suelo, clima y
manejo~\cite{tobler1970}. El rezago de una parcela es el promedio del rendimiento de hasta 10 parcelas de entrenamiento
del mismo régimen a 400\,m o menos, ponderado por el inverso de la distancia. Sin vecinos, recibe la media de su
régimen en el entrenamiento. Elegimos 400\,m porque a 250\,m solo el 25\,\% de las parcelas tiene vecino, mientras que
con radios mayores el promedio mezcla parcelas que ya no se parecen y el error sube.

\textbf{Algoritmo.} Comparamos modelos lineales, de árboles y un proceso gaussiano~\cite{rasmussen2006} con núcleo de
base radial y término de ruido. Sus tres parámetros se ajustan por máxima verosimilitud dentro de cada grupo.

\textbf{Intervalos.} Usamos predicción conformal de Mondrian~\cite{vovk2005,lei2018}, que calibra cada régimen por
separado porque sus errores se distribuyen distinto. A cada predicción le sumamos los cuantiles 0.10 y 0.90 de los
errores fuera de grupo de su régimen.

\textbf{Controles.} Probamos 73 variantes sin corregir por comparaciones múltiples, y elegir la mejor entre muchas
infla su resultado~\cite{cawley2010}. Lo corregimos con validación anidada, que repite la elección entre 12 variantes
dentro de cada grupo externo~\cite{varma2006}. Comparamos modelos con la prueba $t$ corregida de Nadeau y
Bengio~\cite{nadeau2003}, porque las repeticiones comparten parcelas. Una prueba automatizada verifica que alterar las
etiquetas de prueba no cambia las predicciones~\cite{kaufman2012}.

\textbf{Capa económica.} Valuamos la predicción con el precio medio rural del SIAP de 2023 a 2025~\cite{siap2026} y con
los costos del paquete de FIRA de primavera-verano 2022, el más reciente público, actualizados a 2025 con el índice
nacional de precios al consumidor~\cite{agrocostos}.

\section{Resultados y discusión}\label{sec:resultados}

\textbf{Error en validación cruzada.} El modelo obtiene un RMSE de 0.458\,t/ha, un error absoluto medio (MAE) de
0.312\,t/ha y un coeficiente de determinación ($R^2$) de 0.71~\cite{doc_modelo}. Solo con la media de cada régimen el
RMSE es 0.489. La mejora va en la dirección correcta, aunque con 138 parcelas su intervalo de confianza del 95\,\%
todavía incluye el cero (anexo técnico, Tabla~\ref{A-tab:error}).

\textbf{Error esperado en la calificación.} Al repetir toda la elección dentro de la validación el RMSE es 0.490,
prácticamente el de la media por régimen. Por eso esperamos alrededor de 0.49\,t/ha en las 59 parcelas, y el 0.458 es
una cifra optimista. Con solo 59 parcelas el resultado también depende de cuáles toquen, y por azar podría caer entre
0.35 y 0.56~\cite{doc_modelo}.

\textbf{Qué aportó cada pieza.} El régimen es la base. Cuando la validación deja fuera zonas completas, la media por
municipio sube a 0.955\,t/ha y la del régimen se mantiene en 0.489, así que el régimen generaliza y el municipio no. El
rezago espacial mejoró el error en 9 de 10 repeticiones y pasó la prueba de permutación ($p=0.02$). Ningún bloque de
variables satelitales, climáticas o topográficas pasó la regla (Fig.~\ref{A-fig:bloques}). El clima, por ejemplo,
viene en celdas de unos 5\,km y solo distingue zonas, algo que ya hacen el régimen y los vecinos. Entre los algoritmos
comparados, el proceso gaussiano obtuvo el menor error~\cite{doc_modelo}.

\textbf{Cómo verificamos la fuga de información.} Auditamos que todo lo que usa el rendimiento se calcule solo con el
entrenamiento, una prueba automatizada verifica que cambiar las etiquetas de prueba no cambia sus predicciones, y con
los rendimientos revueltos el RMSE no baja de 0.485, contra 0.467 con los reales~\cite{doc_modelo}. Ninguna prueba
descarta la fuga por completo, pero ninguna la encontró.

\textbf{Los intervalos cumplen.} El intervalo del percentil 10 al 90 contiene el rendimiento real en 83\,\% de los
casos, cerca de la meta de 80\,\%~\cite{doc_modelo}. En 18 parcelas, todas del régimen B, la probabilidad de
quedar bajo 2.5\,t/ha, el percentil 10 de ese régimen, supera 10\,\%. La tabla con las 59 predicciones está en el
anexo técnico.

\textbf{El modelo municipal funciona donde no hay datos.} En municipios sin datos del SIAP a menos de 50\,km obtiene un
RMSE de 0.44\,t/ha en promedios municipales, que no se compara con el de parcelas, contra 0.57 del kriging de vecinos,
que no usa clima, y gana en las 10 repeticiones (anexo técnico, Fig.~\ref{A-fig:municipal_modelos}). Como elegimos la
Tmax viendo resultados, repetimos esa elección de forma anidada y la ventaja sigue a favor, con unos 0.07\,t/ha. Cada
grado más de Tmax se asocia con 0.79\,t/ha menos~\cite{doc_generalizable}. El mapa
regional (Fig.~\ref{A-fig:mapas}) muestra que el oriente de Puebla y Tlaxcala, más fresco en verano, tiene el mayor
rendimiento esperado.

\textbf{Por qué el clima no ayudó en las parcelas.} En los nueve municipios del reto la Tmax solo va de 23.3 a
23.9\,$^\circ$C~\cite{doc_generalizable}. A esa escala casi no varía, así que no puede distinguir una parcela de otra.
A escala regional sí.

\textbf{Preguntas guía.} Las bases plantean tres preguntas~\cite{bases}. Como el modelo final no usa índices ni
clima, las respondimos con un bosque aleatorio auxiliar, midiendo cuánto sube su error al revolver cada variable
(anexo técnico, Fig.~\ref{A-fig:indices}). Primero, el índice satelital más importante. Sin el régimen es el NDVI,
porque su historia identifica la zona, y con régimen y vecinos es el índice de sequía NDDI~\cite{gu2007}, aunque apenas
aporta 0.004\,t/ha~\cite{doc_modelo}. Segundo, la variable climática de mayor peso. En las parcelas ninguna aporta una
vez conocido el régimen, y a escala regional es la Tmax de julio y agosto. Tercero, la mejor combinación, que es
régimen más rezago espacial.

\textbf{Qué significa el error para un productor.} Con el precio del SIAP, un error de 0.49\,t/ha equivale a unos
2,700 pesos por hectárea en promedio por parcela, entre 2,100 y 3,500 según el municipio. Es de 7 a 13\,\% del costo de
producción y de 14 a 47\,\% del crédito por hectárea~\cite{repo}. Sirve para ordenar parcelas por riesgo, pero no para
decidir sola sobre una parcela.

\textbf{Aplicación.} En \url{https://agrocebada.alan-pv.com} se consulta el mapa de parcelas, la predicción con su
intervalo y la probabilidad de quedar bajo un umbral que elige el usuario. La aplicación ejecuta el modelo en vivo.

\section{Uso de la ciencia de datos en la inclusión financiera del sector agro}\label{sec:inclusion}

\textbf{De un rendimiento por zona a uno por parcela.} La capa de crédito sustituye el rendimiento esperado del
paquete tecnológico por la predicción de cada parcela con su incertidumbre. Calcula la probabilidad de no cubrir el crédito o el
costo total. La probabilidad combina los errores del modelo con
la fracción de superficie que el municipio pierde históricamente, porque si la cosecha se pierde el ingreso es cero. Un
semáforo la resume en verde bajo 10\,\%, amarillo de 10 a 30\,\% y rojo desde 30\,\%. Con el crédito, 23 de las 59
parcelas salen en verde, 30 en amarillo y 6 en rojo. Con el costo total salen 55 en rojo, porque el paquete incluye
conceptos no financiados, como la renta.

\textbf{A quién sirve.} La comercializadora que compra el grano puede priorizar visitas de campo. El analista ordena
parcelas por riesgo y el productor ve cuánto debe cosechar para pagar.

\textbf{Alcance.} La probabilidad describe la incertidumbre de un año parecido a 2025 y no incluye sequías, heladas ni
cambios de precio. Por eso no es una probabilidad de incumplimiento ni sirve para tarificar seguros, que además
enfrentan riesgo base~\cite{jensen2017}.

\textbf{Momento de uso.} El modelo usa rendimientos de vecinas de la misma temporada, que existen solo después de su
cosecha. Es una estimación de cierre de temporada. Para usarlo antes de cosechar harían falta rendimientos de ciclos
anteriores, que el reto no incluye.

\textbf{Escala municipal.} Con el SIAP de 2003 a 2025 ajustamos un modelo municipal~\cite{siap2026}. En
municipios a más de 50\,km de los de entrenamiento, la temperatura máxima de julio y agosto más kriging de residuos
obtuvo un RMSE de 0.44\,t/ha contra 0.57 del kriging solo. Al repetir de forma anidada
la elección de esa variable, la ventaja baja a unos 0.07\,t/ha.

\section{Conclusiones}\label{sec:conclusiones}

Construimos un modelo sencillo para estimar el rendimiento de 59 parcelas de cebada de temporal. Usa dos datos, el
régimen de registro del municipio y el rendimiento de las parcelas vecinas, combinados con un proceso gaussiano.
Esperamos un RMSE de 0.46 a 0.49\,t/ha en las 59 parcelas.

Cada pieza entró porque mejoró la validación de forma consistente, y cada idea que quedó fuera lo hizo porque no la
mejoró. Así quedaron fuera las variables satelitales y climáticas, que no aportaron una vez conocidos el régimen y los
vecinos.

Nuestros resultados sugieren que el límite de este reto está más en las etiquetas que en el algoritmo. El modelo mejora
poco a la media por régimen, y el redondeo y la fuente mixta del rendimiento pueden agregar un ruido que no pudimos
medir. Conocer la fuente de cada registro, la fecha de siembra o tener cortes de cosecha medidos permitiría comprobarlo,
y podría ayudar más que un modelo más complejo.

Sobre las preguntas guía, el NDVI y la lluvia pesan solo cuando falta el régimen, porque indican la zona. Con el régimen
y los vecinos, el NDDI es el índice que más aporta, pero su efecto es muy pequeño. La mejor combinación es régimen más
rezago espacial.

Para un productor, el error equivale a unos 2,500 pesos por hectárea. La estimación sirve para ordenar parcelas y
orientar la revisión de un crédito, junto con el juicio del analista.

Como trabajo futuro proponemos validar con el rendimiento real de las 59 parcelas cuando FIRA lo publique, sumar más
temporadas para estimar antes de la cosecha y probar la capa de crédito con datos reales.


% Generado por plataforma.py. Estilo IEEE, en orden de primera cita.
\begin{thebibliography}{99}

\bibitem{bases}
FIRA, \emph{Reto AgroCebada 2026. Bases de participación y descripción del conjunto de datos},
  2026.

\bibitem{sader2025}
Secretaría de Agricultura y Desarrollo Rural, ``Del campo a la copa: el recorrido de la cebada,'' \emph{gob.mx}, 11 de noviembre de 2025. [En línea]. Disponible: \url{https://www.gob.mx/agricultura/articulos/del-campo-a-la-copa-el-recorrido-de-la-cebada}

\bibitem{preguntas}
FIRA, \emph{Sesiones de preguntas y respuestas del Reto AgroCebada 2026 y correos de seguimiento a los equipos participantes}, septiembre y octubre de 2026 (último correo del 14 de octubre).

\bibitem{drusch2012}
M. Drusch \emph{et al.}, ``Sentinel-2: ESA's Optical High-Resolution Mission for GMES Operational Services,'' \emph{Remote Sensing of Environment}, vol.~120, pp.~25--36, 2012.

\bibitem{frazier2021}
A.~E. Frazier y B.~L. Hemingway, ``A Technical Review of Planet Smallsat Data: Practical Considerations for Processing and Using PlanetScope Imagery,'' \emph{Remote Sensing}, vol.~13, no.~19, p.~3930, 2021.

\bibitem{funk2015}
C. Funk \emph{et al.}, ``The climate hazards infrared precipitation with stations---a new environmental record for monitoring extremes,'' \emph{Scientific Data}, vol.~2, art.~150066, 2015.

\bibitem{verdin2020}
A. Verdin \emph{et al.}, ``Development and validation of the CHIRTS-daily quasi-global high-resolution daily temperature data set,'' \emph{Scientific Data}, vol.~7, no.~1, art.~303, 2020.

\bibitem{vanklompenburg2020}
T. van Klompenburg, A. Kassahun y C. Catal, ``Crop yield prediction using machine learning: A systematic literature review,'' \emph{Computers and Electronics in Agriculture}, vol.~177, p.~105709, 2020.

\bibitem{weiss2020}
M. Weiss, F. Jacob y G. Duveiller, ``Remote sensing for agricultural applications: A meta-review,'' \emph{Remote Sensing of Environment}, vol.~236, p.~111402, 2020.

\bibitem{lobell2020}
D.~B. Lobell \emph{et al.}, ``Eyes in the Sky, Boots on the Ground: Assessing Satellite‐ and Ground‐Based Approaches to Crop Yield Measurement and Analysis,'' \emph{American Journal of Agricultural Economics}, vol.~102, no.~1, pp.~202--219, 2020.

\bibitem{agrocostos}
Fideicomisos Instituidos en Relación con la Agricultura, ``Agrocostos.'' [En línea]. Disponible: \url{https://www.fira.gob.mx/Nd/Agrocostos.jsp}. Consultado el 15 de octubre de 2026.

\bibitem{ecuaciones}
FIRA, \emph{Ecuaciones de índices satelitales. Reto AgroCebada FIRA 2026}, 2026.

\bibitem{zadoks1974}
J.~C. Zadoks, T.~T. Chang y C.~F. Konzak, ``A decimal code for the growth stages of cereals,'' \emph{Weed Research}, vol.~14, no.~6, pp.~415--421, 1974.

\bibitem{roberts2017}
D.~R. Roberts \emph{et al.}, ``Cross-validation strategies for data with temporal, spatial, hierarchical, or
phylogenetic structure,'' \emph{Ecography}, vol.~40, no.~8, pp.~913--929, 2017.

\bibitem{ploton2020}
P. Ploton \emph{et al.}, ``Spatial validation reveals poor predictive performance of large-scale ecological mapping models,'' \emph{Nature Communications}, vol.~11, no.~1, art.~4540, 2020.

\bibitem{ojala2010}
M.~Ojala y G.~C. Garriga, ``Permutation tests for studying classifier performance,'' \emph{Journal of Machine
Learning Research}, vol.~11, pp.~1833--1863, 2010.

\bibitem{tobler1970}
W.~R. Tobler, ``A computer movie simulating urban growth in the Detroit region,'' \emph{Economic Geography},
vol.~46, pp.~234--240, 1970.

\bibitem{rasmussen2006}
C.~E. Rasmussen y C.~K.~I. Williams, \emph{Gaussian Processes for Machine Learning}. Cambridge, MA, EE.~UU.: MIT
Press, 2006.

\bibitem{vovk2005}
V.~Vovk, A.~Gammerman y G.~Shafer, \emph{Algorithmic Learning in a Random World}. Nueva York, NY, EE.~UU.:
Springer, 2005.

\bibitem{lei2018}
J. Lei \emph{et al.}, ``Distribution-Free Predictive Inference for Regression,'' \emph{Journal of the American Statistical Association}, vol.~113, no.~523, pp.~1094--1111, 2018.

\bibitem{cawley2010}
G.~C. Cawley y N.~L.~C. Talbot, ``On over-fitting in model selection and subsequent selection bias in performance
evaluation,'' \emph{Journal of Machine Learning Research}, vol.~11, pp.~2079--2107, 2010.

\bibitem{varma2006}
S.~Varma y R.~Simon, ``Bias in error estimation when using cross-validation for model selection,'' \emph{BMC
Bioinformatics}, vol.~7, art.~91, 2006.

\bibitem{nadeau2003}
C.~Nadeau e Y.~Bengio, ``Inference for the generalization error,'' \emph{Machine Learning}, vol.~52, no.~3,
pp.~239--281, 2003.

\bibitem{kaufman2012}
S.~Kaufman, S.~Rosset, C.~Perlich y O.~Stitelman, ``Leakage in data mining: Formulation, detection, and
avoidance,'' \emph{ACM Transactions on Knowledge Discovery from Data}, vol.~6, no.~4, art.~15, 2012.

\bibitem{siap2026}
SIAP, ``Producción agrícola municipal, cierre 2025,'' Datos abiertos, 2026. Consultado el 4 de
  octubre de 2026. Disponible en \url{https://nube.agricultura.gob.mx/datosAbiertos/Agricola.php}

\bibitem{gu2007}
Y. Gu \emph{et al.}, ``A five‐year analysis of MODIS NDVI and NDWI for grassland drought assessment over the central Great Plains of the United States,'' \emph{Geophysical Research Letters}, vol.~34, no.~6, art.~2006GL029127, 2007.

\bibitem{jensen2017}
N. Jensen y C. Barrett, ``Agricultural Index Insurance for Development,'' \emph{Applied Economic Perspectives and Policy}, vol.~39, no.~2, pp.~199--219, 2017.

\end{thebibliography}


\end{document}
